% !TEX TS-program = XeLaTeX
% use the following command:
% all document files must be coded in UTF-8
\documentclass[portuguese]{textolivre}
% build HTML with: make4ht -e build.lua -c textolivre.cfg -x -u article "fn-in,svg,pic-align"

\journalname{Texto Livre}
\thevolume{19}
%\thenumber{1} % old template
\theyear{2026}
\receiveddate{\DTMdisplaydate{2025}{11}{30}{-1}} % YYYY MM DD
\accepteddate{\DTMdisplaydate{2026}{1}{19}{-1}}
\publisheddate{\DTMdisplaydate{2026}{7}{22}{-1}}
\corrauthor{Leandro Paz da Silva}
\articledoi{10.1590/1983-3652.2026.63156}
%\articleid{NNNN} % if the article ID is not the last 5 numbers of its DOI, provide it using \articleid{} commmand 
% list of available sesscions in the journal: articles, dossier, reports, essays, reviews, interviews, editorial
\articlesessionname{dossier}
\runningauthor{Silva, Villarroel e Okuyama} 
%\editorname{Leonardo Araújo} % old template
\sectioneditorname{Zaira Bomfante dos Santos~\orcid{0000-0002-6162-8489}}
\layouteditorname{Leonardo Araujo~\orcid{0000-0003-3884-2177}}

\title{INATICON: o uso de agentes conversacionais com inteligência artificial para potencializar a docência no ensino superior atuando sobre multiletramentos}
\othertitle{INATICON: the use of conversational agents with artificial intelligence to enhance teaching in higher education acting on multiliteracies}
% if there is a third language title, add here:
%\othertitle{Artikelvorlage zur Einreichung beim Texto Livre Journal}

\author[1]{Leandro Paz da Silva~\orcid{0000-0002-4979-3625}\thanks{Email: \href{mailto:l.silva@ufrgs.br}{l.silva@ufrgs.br}}}
\author[2]{Márcia Amaral Corrêa Ughini Villarroel~\orcid{0000-0002-6676-3698}\thanks{Email: \href{mailto:marcia.correa@sertao.ifrs.edu.br}{marcia.correa@sertao.ifrs.edu.br}}}
\author[1]{Fabio Yoshimitsu Okuyama~\orcid{0000-0001-7427-2525}\thanks{Email: \href{mailto:fabio.okuyama@poa.ifrs.edu.br}{fabio.okuyama@poa.ifrs.edu.br}}}
\affil[1]{Instituto Federal de Educação, Ciência e Tecnologia do Rio Grande do Sul, Porto Alegre, RS, Brasil.}
\affil[2]{Instituto Federal de Educação, Ciência e Tecnologia do Rio Grande do Sul, Sertão, RS, Brasil.}

\addbibresource{article.bib}
% use biber instead of bibtex
% $ biber article

% used to create dummy text for the template file
\definecolor{dark-gray}{gray}{0.35} % color used to display dummy texts
\usepackage{lipsum}
\SetLipsumParListSurrounders{\colorlet{oldcolor}{.}\color{dark-gray}}{\color{oldcolor}}

% used here only to provide the XeLaTeX and BibTeX logos
\usepackage{hologo}

% if you use multirows in a table, include the multirow package
\usepackage{multirow}

% provides sidewaysfigure environment
\usepackage{rotating}

% CUSTOM EPIGRAPH - BEGIN 
%%% https://tex.stackexchange.com/questions/193178/specific-epigraph-style
\usepackage{epigraph}
\renewcommand\textflush{flushright}
\makeatletter
\newlength\epitextskip
\pretocmd{\@epitext}{\em}{}{}
\apptocmd{\@epitext}{\em}{}{}
\patchcmd{\epigraph}{\@epitext{#1}\\}{\@epitext{#1}\\[\epitextskip]}{}{}
\makeatother
\setlength\epigraphrule{0pt}
\setlength\epitextskip{0.5ex}
\setlength\epigraphwidth{.7\textwidth}
% CUSTOM EPIGRAPH - END

% to use IPA symbols in unicode add
%\usepackage{fontspec}
%\newfontfamily\ipafont{CMU Serif}
%\newcommand{\ipa}[1]{{\ipafont #1}}
% and in the text you may use the \ipa{...} command passing the symbols in unicode

% LANGUAGE - BEGIN
% ARABIC
% for languages that use special fonts, you must provide the typeface that will be used
% \setotherlanguage{arabic}
% \newfontfamily\arabicfont[Script=Arabic]{Amiri}
% \newfontfamily\arabicfontsf[Script=Arabic]{Amiri}
% \newfontfamily\arabicfonttt[Script=Arabic]{Amiri}
%
% in the article, to add arabic text use: \textlang{arabic}{ ... }
%
% RUSSIAN
% for russian text we also need to define fonts with support for Cyrillic script
% \usepackage{fontspec}
% \setotherlanguage{russian}
% \newfontfamily\cyrillicfont{Times New Roman}
% \newfontfamily\cyrillicfontsf{Times New Roman}[Script=Cyrillic]
% \newfontfamily\cyrillicfonttt{Times New Roman}[Script=Cyrillic]
%
% in the text use \begin{russian} ... \end{russian}
% LANGUAGE - END

% EMOJIS - BEGIN
% to use emoticons in your manuscript
% https://stackoverflow.com/questions/190145/how-to-insert-emoticons-in-latex/57076064
% using font Symbola, which has full support
% the font may be downloaded at:
% https://dn-works.com/ufas/
% add to preamble:
% \newfontfamily\Symbola{Symbola}
% in the text use:
% {\Symbola }
% EMOJIS - END

% LABEL REFERENCE TO DESCRIPTIVE LIST - BEGIN
% reference itens in a descriptive list using their labels instead of numbers
% insert the code below in the preambule:
%\makeatletter
%\let\orgdescriptionlabel\descriptionlabel
%\renewcommand*{\descriptionlabel}[1]{%
%  \let\orglabel\label
%  \let\label\@gobble
%  \phantomsection
%  \edef\@currentlabel{#1\unskip}%
%  \let\label\orglabel
%  \orgdescriptionlabel{#1}%
%}
%\makeatother
%
% in your document, use as illustraded here:
%\begin{description}
%  \item[first\label{itm1}] this is only an example;
%  % ...  add more items
%\end{description}
% LABEL REFERENCE TO DESCRIPTIVE LIST - END


% add line numbers for submission
%\usepackage{lineno}
%\linenumbers

\begin{document}
\maketitle

\begin{polyabstract}
\begin{abstract}
O presente estudo investiga o potencial de agentes conversacionais com Inteligência Artificial (IA), como ferramenta de auxílio à docência no ensino superior. A pesquisa, fundamentada em Design Science Research (DSR), resultou no desenvolvimento do INATICON, um produto técnico que media a interação docente com a IA. A partir de intervenções com professores da Rede Federal de Ensino, a análise dos dados, realizada sob a ótica da Análise Crítica do Discurso (ACD), revelou como resultado principal um fenômeno análogo ao efeito Dunning-Kruger, no qual os participantes, inicialmente superconfiantes, passaram a uma postura de "consciência crítica" após a interação com a ferramenta. Os resultados indicam que o uso mediado da IA, focado na potência ontológica e na andragogia, pode fomentar uma prática docente mais reflexiva e eficiente, promovendo o uso ético e responsável da tecnologia na educação.

\keywords{Agentes conversacionais \sep Inteligência Artificial \sep Docência \sep Andragogia \sep Multiletramentos}
\end{abstract}

\begin{english}
\begin{abstract}
This study investigates the potential of conversational agents with Artificial Intelligence (AI) to support teaching in higher education. The research, based on Design Science Research (DSR), resulted in the development of INATICON, a technical product that mediates teacher interaction with AI. From interventions with professors of the Brazilian Federal Education Network, data analysis, conducted from the perspective of Critical Discourse Analysis (CDA), revealed as a main result a phenomenon analogous to the Dunning-Kruger effect, in which participants, initially overconfident, shifted to a stance of "critical awareness" after interacting with the tool. The results indicate that the mediated use of AI, focused on ontohological potency and andragogy, can foster a more reflective and efficient teaching practice, promoting the ethical and responsible use of technology in education.

\keywords{Conversational agents \sep Artificial Intelligence \sep Teaching \sep Andragogy \sep Multiliteracies}
\end{abstract}
\end{english}
% if there is another abstract, insert it here using the same scheme
\end{polyabstract}

\section{Introdução}\label{sec-intro}
A emergência dos agentes conversacionais (também conhecidos como \textit{chatbots}) com IA nos anos 2020, como o ChatGPT da OpenAI, impulsionados por modelos robustos de IA baseados em transformadores textuais generativos \cite{bochie2020aprendizado}, desencadeou um amplo debate e, por vezes, uma "histeria" no campo educacional \cite{selwyn2024limits}. O acesso popularizado a essas ferramentas gerou questionamentos constantes sobre seu impacto na docência, na autoria e na qualidade do discurso (seja do estudante ou do professor), expondo a carência de trabalhos nacionais ''baseados em evidências" que ofereçam uma resposta "austera e sólida" para o uso planejado dessas tecnologias nas Instituições Federais de Ensino Superior. Dessa forma, o presente estudo de caráter interdisciplinar, como parte do Mestrado Profissional de Informática na Educação do Instituto Federal de Educação, Ciência e Tecnologia (IFRS), que visa a criação de um produto técnico – busca responder à questão fundamental: como o uso de agentes conversacionais com IA pode potencializar a atuação docente no ensino superior no contexto atual?

\section{Referencial teórico e revisão da literatura}\label{sec-normas}
A fundamentação teórica desta pesquisa ancora-se na ACD que transcende o estatuto de mera técnica ou recurso metodológico \cite{silva2017analise} para se estabelecer como uma abordagem teórico-metodológica com profundo compromisso político e epistemológico \cite{mesquita2020analise}. A ACD insere-se no quadro das teorias sociais críticas, fundamentada na concepção da linguagem como prática social \cite{fairclough1989language,fairclough2001discourse}. O objetivo central dessa vertente é o estudo do discurso como um momento da vida social, visando ao desvelamento das relações de poder, dominação e desigualdade que são produzidas ou reproduzidas por meio de textos e falas \cite{vandijk2001critical}. A ACD atua como uma lente que conecta o texto à realidade social \cite{fairclough2006language}, investigando como as assimetrias de poder e as ideologias são manifestadas no discurso \cite{chouliaraki1999discourse}. Essa perspectiva é fundamental para o exame dos letramentos e multiletramentos \cite{rojo2012pedagogia}, que são centrais para analisar as interações textuais entre docentes e agentes conversacionais de Inteligência Artificial (IA). O foco no Letramento Crítico \cite{janks2010literacy,janks2012importance} permite refletir sobre as implicações da IA na autoria e na prática docente, dada a natureza ideológica e constitutiva do social que a ACD busca desnaturalizar \cite{wodak2009methods}.

A coerência científica e a relevância do objeto de estudo são sustentadas por uma Revisão Sistemática de Literatura (RSL) na área de Informática na Educação, concluída em 2024. Os resultados da RSL revelaram uma lacuna significativa na produção científica. Foi constatada uma baixa apropriação aos agentes conversacionais por parte dos docentes para fins de ensino no contexto do Ensino Superior. O panorama obtido sugere que o volume de trabalhos é maior quando concentrado no aprender com agentes de IA em comparação com o foco no ensinar com IA \cite{silveira2019uso}. Além disso, a RSL notou uma baixa taxa de citação ou referenciação entre os trabalhos na área, o que sugere que os estudos existentes tendem a ser isolados e multifacetados. Muitos trabalhos são da área de computação sobre aprendizado de máquina, com pouca relação direta com a prática pedagógica ou o trabalho docente \cite{zhai2021review}. Tais achados reforçam a demanda científica por uma pesquisa organizada e focada na IA como meio para potencialização da docência no Ensino Superior \cite{guimaraes2023metodologias}, justificando assim a necessidade de preencher a lacuna de apropriação crítica e organizada por parte dos educadores.

\subsection{O conceito de potência e a ação docente}\label{sec-conduta}
O estudo fundamenta-se na filosofia de Spinoza, concebendo a potência (\textit{conatus}, em latim) como a própria "ser-ação" do ente. Nessa perspectiva, \textit{conatus} é o esforço contínuo de cada ser em perseverar na existência e atualizar suas capacidades, o que se traduz em um modelo de liberdade como a "plena realização da própria natureza" \cite{spinoza2007etica}. Em vez de focar na moralidade ou legalidade do uso das IAs (o "devemos ou não usar?"), a pesquisa concentra-se em como o uso estratégico dessas ferramentas pode ampliar a liberdade e a eficácia do professor, ou seja, a potência docente, tornando o indivíduo a "causa adequada" de sua prática \cite{barros2018revisitando,spinoza2007etica}.


\subsection{Experiência e andragogia}\label{sec-fmt-manuscrito}
O recorte metodológico e o foco no público-alvo são sustentados pela perspectiva andragógica, uma vez que os participantes são principalmente docentes adultos das Instituições Federais de Ensino Superior. A andragogia (Knowles \textit{apud} \textcite{barros2018revisitando}) reconhece a experiência acumulada do adulto como um recurso fundamental, que, filosoficamente, pode ser vista como um aumento da potência de agir. Ademais, a pesquisa alinha-se ao pragmatismo de John Dewey, entendendo o ensino como um "chamado para a experiência" e a "prática social" \cite{dewey1959democracia,dewey1979experiencia}, o que converge com as metodologias ativas que valorizam o processo experimental de "tentar e errar", essenciais ao interagir com agentes conversacionais \cite{bacich2018metodologias}.


\subsection{Letramentos, multiletramentos e dialogismo}\label{sec-formato}
As reflexões sobre a interação textual e o discurso com a IA são ancoradas nos conceitos de multiletramentos \cite{rojo2012pedagogia} e dialogismo \cite{bakhtin2016generos}. Agentes conversacionais com IA têm o potencial de ampliar os multiletramentos, especialmente quando combinam texto escrito (grafado) com a geração de imagens, como observado em plataformas de IA generativas e altamente criativas, apresentando uma grande margem para o fazer lúdico e o fazer criativo enquanto ludoletramento \cite{zagal2010ludoliteracy} ou domínio semiótico amplo \cite{gee2004situated}. O uso do dialogismo, em particular o conceito de dialogismo mais ativo na conversação \cite{suleman2016negotiation}, serve para investigar como o agente conversacional pode fomentar processos reflexivos e autoavaliativos ao fornecer feedback imediato, atuando como um "mediador ou dialogador" na construção do conhecimento pelo usuário \cite{costa2009proposta,junior2017ontologias}.


\section{Metodologia}\label{sec-modelo}
A metodologia de pesquisa estrutura-se de forma mista e complementar, tendo como eixo central a DSR, que é a abordagem principal para a construção e validação do produto técnico INATICON – nesse sentido, a DSR é um método de natureza prescritiva que será mantido de forma cíclica e constante, desde a intervenção inicial até a validação final do artefato \cite{pimentel2020design,simon1988science}. Para o rigor do tratamento e interpretação dos dados textuais gerados, será empregada a ACD, que não é vista apenas como uma técnica, mas como uma abordagem teórico-metodológica crítica \cite{mesquita2020analise,queiroz2014analise}. A ACD atuará como lente principal para a análise da apropriação crítica do discurso por parte dos docentes, focando nas questões de letramentos e multiletramentos \cite{janks2010literacy,rojo2012pedagogia}. Adicionalmente, será utilizada a Mineração de Textos \cite{reategui2017apoio} como metodologia secundária e complementar, essencial para a organização e análise em panorama dos dados textuais provenientes das interações dos participantes com os agentes conversacionais, conferindo um suplemento quantificável à microanálise linguística da ACD.

\section{Metodologia de pesquisa e desenvolvimento – produto técnico}\label{sec-organizacao}
Para garantir o rigor metodológico na construção e validação do artefato INATICON, a pesquisa adota a DSR como núcleo metodológico \cite{pimentel2020design,simon1988science}. A DSR foca na elaboração de um produto técnico (INATICON) que é continuamente aprimorado em ciclos a partir das interações e dos retornos dos participantes da pesquisa. O artefato técnico resultante é de caráter prescritivo e adaptável, visando orientar o uso planejado de agentes conversacionais com IA para fins de docência. Em síntese, o artefato INATICON seria aqui um mediador para usar um \textit{chat} com IA como mediador educacional, em outras palavras. Ao longo do desenvolvimento da pesquisa o produto técnico foi transformado em um produto para ser utilizado dentro de ambientes como o Moodle, por meio do pacote H5P, sendo um .h5p aberto (\Cref{fig1}).

\begin{figure}[htbp]
\centering
\begin{minipage}{.9\textwidth}
 \includegraphics[width=\textwidth]{Fig1.png}
 \caption{Produto Técnico INATICON em desenvolvimento.}
 \label{fig1}
 \source{Moodle IFRS POA/ Autor.}
\end{minipage}
\end{figure}

\subsection{Metodologia de análise de dados}\label{sec-organizacao-latex}
Em paralelo à DSR, para analisar as interações textuais (\textit{logs} de \textit{chat}) e o desenvolvimento da consciência crítica dos docentes, a pesquisa emprega a (ACD) aplicada ao ensino \cite{janks2010literacy,janks2012importance,queiroz2014analise}. A ACD investiga como as entradas (\textit{prompts}) e as saídas (respostas da IA) operam no sentido e no discurso dos seus usuários, examinando a apropriação crítica dos multiletramentos. Essa análise é suplementada pela Mineração de Textos \cite{reategui2017apoio}, uma técnica que organiza e analisa dados textuais para solidificar a análise linguística com base em dados, relevante tanto para ciências duras quanto para humanidades.

Considerando a ACD em práxis, a análise das interações entre os docentes participantes da pesquisa e as plataformas de IA generativa revela questões cruciais de poder e linguagem, enquanto metodologia que lida diretamente com as dinâmicas de poder textuais e discursivas. As plataformas de IA sendo modelos robustos (como o GPT) e, frequentemente, proprietários, atuam como fontes de autoridade discursiva que os docentes precisam testar e confrontar no caso de nossa demanda de pesquisa. A ACD é empregada para examinar como as entradas e saídas de texto operam no sentido e no discurso dos seus usuários, e como essa interação afetaria também os multiletramentos envolvidos. Por exemplo, o teste realizado pela participante Anônima A, que delegou a criação de um plano de ensino de 18 semanas, constituiu uma crítica avaliativa destinada a testar a competência epistêmica e a autoridade da IA. Em contraste, a participante Anônima C, ao usar o recurso "gameficador", expôs (na sua percepção) a "ideologia" da ferramenta – que tende a gerar soluções teóricas desvinculadas da práxis docente e das limitações contextuais, como a falta de consideração do tempo de aula. A dinâmica de poder, portanto, é disputada na capacidade do docente de usar o \textit{prompt} como uma ferramenta de investigação crítica, revelando as lacunas funcionais do discurso algorítmico.

O conflito dialético central reside na tensão entre a aceleração das IAs, coordenadas pelas \textit{Big Techs} no capitalismo do século XXI, e a demanda por uma integração pedagógica, que pretendemos que seja refletida e "austera" seguindo as diretrizes da \textcite{unesco2019beijing,unesco2023chatgpt,unesco2024guia}. Nesse sentido, a pesquisa emerge como uma resposta à "histeria" do boom da IA generativa transformacional, também conhecida por \textit{Large Language Model}, modeladas em \textit{Machine Learning} ou \textit{Deep Learning} principalmente, buscando conferir benefícios ao "bem público" e evitar a dominação obscura lucrativa das \textit{Big Techs} sobre o trabalho docente e a sociedade como um todo.

Já o conflito dialógico \cite{bakhtin2016generos} manifesta-se na luta pela autoria e pela integração contextual em um ambiente que deveria promover o dialogismo, mas que se baseia em modelos proprietários. A despeito do potencial dialógico da IA para promover resposta imediata e processos autoavaliativos, o desafio está em sua natureza: o sistema falha ao gerar soluções teóricas sem perguntar ou internalizar o contexto real do professor. A crítica funcional da Anônima C ressaltou que a IA é incapaz de perguntar quanto tempo a docente teria disponível, expondo uma lacuna crucial entre a saída textual da IA e a prática de sala de aula. Para que a ferramenta possa de fato "potencializar a atuação docente", ela deve ser concebida como um meio que respeite a autonomia e a experiência acumulada do docente adulto (andragogia), o que requer que o professor atue ativamente para ensinar à IA o seu contexto antes que ela possa servir como um auxiliar de design instrucional.


\subsection{Intervenções para a coleta de dados}\label{sec-titulo}
O processo de validação envolveu intervenções em formato de oficinas com docentes da Instituições Federais de Ensino Superior. Durante essas oficinas, os dados foram coletados por meio de questionários após o uso do artefato INATICON e \textit{logs} de \textit{chat} com IA permitindo a análise das interações textuais pela ACD (enquanto metodologia central de análise) e a Mineração de Textos (enquanto método acessório complementar). Esta metodologia rigorosa é fundamental para superar a análise superficial, e, garantindo a reprodutibilidade do experimento com o artefato posteriormente.

\subsection{Considerações éticas}\label{sec-autores}
A pesquisa, que envolve sujeitos humanos (andragógicos), foi devidamente aprovada pelo Comitê de Ética em Pesquisa da instituição do autor. O trabalho respeita a Lei Geral de Proteção de Dados Pessoais (LGPD) \cite{brasil2018lgpd}, garantindo que os dados em texto dos sujeitos sejam individuais e anonimizados, utilizando pseudônimos. Além disso, o autor comprometeu-se a seguir as recomendações da Texto Livre sobre o uso de ferramentas de IA informando e assumindo a responsabilidade pelo conteúdo gerado, uma vez que "apenas humanos podem ser autores".

\section{O artefato INATICON: design e recursos}\label{sec-idioma}
O INATICON é um artefato desenvolvido em DSR, um produto técnico, adaptável e replicável, projetado para ser um instrumento mediador que fornece uma orientação técnica e interdisciplinar aos docentes. O nome, que significa "contra o inato" (focado no artifício e na potência), reflete a projeção estética que busca inspiração no século XVII, um período de grande florescimento nas artes e ciências, para tratar o boom das IAs do século XXI. Seu design instrucional é baseado no modelo "contextualizado" \cite{filatro2010design}, sendo lúdico e seguindo parâmetros de gamificação \cite{carvalho2020livro}. Os recursos propostos, detalhados em fluxogramas, são extensões das diretrizes do GIR da \textcite{unesco2023chatgpt} e incluem recursos como o Co-designer (para assistência em planejamento didático), o Ativador (para aproximar o planejamento de metodologias ativas) e o Motivador (para gerar atividades que aumentem o engajamento do aluno com desafios).



\section{Validação, resultados e discussão}\label{sec-resumo}
Os resultados da pesquisa aplicada, conduzida por meio de oficinas com docentes, revelaram uma dificuldade de percepção das próprias habilidades na conversação com a IA, um fenômeno análogo ao efeito Dunning-Kruger. Este viés cognitivo, no qual indivíduos com competência limitada em uma área superestimam suas habilidades, mostrou-se um achado central, demonstrando tanto as fragilidades quanto a eficácia do uso da ferramenta.

O efeito Dunning-Kruger \cite{kruger1999unskilled} postula que a falta de conhecimento em um domínio específico impede o indivíduo de reconhecer a própria incompetência. Essa "dupla maldição" – ausência de conhecimento e de consciência sobre essa ausência – resulta em uma autoconfiança inflada e desproporcional à habilidade real. Na pesquisa qualitativa, tal efeito pode impactar tanto o pesquisador, que pode superestimar sua capacidade de análise, quanto os participantes, que podem se apresentar como mais conhecedores do que são, introduzindo vieses nos dados \cite{longweni2025managerial}.

A análise dos dados coletados demonstrou essa dinâmica de forma clara. Por um lado, a própria IA apresenta dificuldade de autopercepção em um \textit{chat}, o que a leva a "alucinar" e a ter dificuldade em manter uma conversa com dados confiáveis. Por outro lado, os participantes da pesquisa – em sua maioria docentes que não utilizavam tecnologias com IA – iniciaram a interação com altas expectativas e confiança na ferramenta. Contudo, a análise das conversas revelou uma queda consistente na confiança dos participantes no uso da IA após a oficina. Essa queda é interpretada não como desestímulo, mas como o início do pensamento crítico, movendo os docentes de uma "confiança ingênua" para uma "consciência crítica", um indicador do sucesso pedagógico da intervenção.

A análise dos logs de \textit{chat} corroborou este engajamento complexo. Para sistematizar a análise, os excertos foram selecionados com base na sua representatividade em ilustrar as diferentes posturas críticas dos docentes ao interagir com a IA. Foram identificados três tipos de crítica:

\begin{enumerate}
    \item \textit{Crítica avaliativa}: Focada em testar a competência da IA em tarefas complexas. Exemplo: a docente Anônima A, que delegou à IA a criação de um plano de ensino de 18 semanas para avaliar sua capacidade epistêmica.
    \item \textit{Crítica utilitária}: Centrada na valorização da automação e da velocidade. Exemplo: o docente Anônimo B, que utilizou a ferramenta para otimizar tarefas repetitivas.
    \item \textit{Crítica funcional}: Destinada a expor a incapacidade da IA de considerar o contexto real da prática docente. Exemplo: a docente Anônima C, que, ao usar o recurso "gameficador", notou a ausência de consideração sobre o tempo de aula disponível, revelando a lacuna entre a solução teórica da IA e a realidade da sala de aula.
\end{enumerate}

Em outras palavras, a análise dos formulários revelou a manifestação recorrente do fenômeno: no início da oficina, participantes com baixo conhecimento sobre os tópicos abordados apresentaram níveis notavelmente elevados de autoconfiança, que se inverteram ao final da intervenção. Contudo, a simples identificação desse viés não é suficiente. Utilizando a ACD, é possível examinar como essas dinâmicas se materializam linguisticamente. A ACD permite "desmascarar" a construção discursiva da autoridade baseada na confiança, e não na competência, e seu papel nas estruturas de poder, analisando os dados da pesquisa qualitativa.

Para \textcite{vandijk2008discurso}, as ideologias são representações sociais compartilhadas que se manifestam no discurso e legitimam a dominação. A ACD, nesse sentido, busca tornar explícitas as estruturas de poder e as ideologias subjacentes ao texto e à fala. A integração da ACD com os conceitos psicológicos permite uma análise aprofundada de como o viés análogo ao efeito \textit{Dunning-Kruger} se perpetua no âmbito social. A ACD pode identificar as estratégias discursivas utilizadas por indivíduos superconfiantes, mas pouco informados, como o uso de simplificações excessivas, linguagem categórica ("é óbvio que...") e a apropriação de jargões técnicos para construir uma "fachada" de autoridade.

Conforme postulado por \textcite{vandijk2008discurso}, o controle do discurso está intrinsecamente ligado ao poder. A confiança pode se tornar um capital social que garante acesso a plataformas de discussão, silenciando vozes mais competentes, porém menos assertivas. Para \textcite{fairclough2003analysing}, o discurso é central para a reprodução de relações sociais e de poder. Ao analisar criticamente o viés observado, a ACD revela como a valorização social da assertividade infundada sobre a competência ponderada contribui para a manutenção de estruturas de poder que recompensam a performance em detrimento do conhecimento substancial.

Portanto, o fenômeno observado não é apenas uma curiosidade psicológica; ele possui consequências materiais e discursivas profundas. A ACD fornece um caminho crítico para a compreensão de como a autoconfiança infundada se traduz em poder social. Ao desvelar as estratégias linguísticas e ideológicas subjacentes, a ACD permite uma leitura mais crítica da realidade, expondo os mecanismos pelos quais a competência real é, muitas vezes, silenciada pelo ruído da certeza ignorante.

Por fim, a jornada com o INATICON também pode levar a não recomendação do uso de tecnologias com IA por parte de docentes mais próximos de vieses instrumentais ou "acelerados". Na prática, o INATICON sugere ao docente que é possível "desacelerar" caso se sinta inseguro, ou que exercite o \textit{chat} com IA sem grandes pretensões, superando uma teleologia imediata. Talvez não recomendar seja uma forma de repensar a busca por textos e imagens, em uma missão de se munir de outros letramentos, para além das tecnologias com IA.

\section{Conclusão}\label{sec-secoes}
O presente trabalho propõe o INATICON, um instrumento mediador validado pela DSR e analisado pela ACD, que busca potencializar a atuação docente no Ensino Superior. Atualmente, o artefato encontra-se em fase avançada de desenvolvimento, sendo testado como um recurso .h5p no ambiente Moodle da instituição. Os resultados da pesquisa, com destaque para a identificação de um viés análogo ao Efeito Dunning-Kruger, demonstram que a investigação sobre a percepção docente frente ao aceleracionismo das IAs contribui para uma "desaceleração e racionalização crítica do tema" \cite{selwyn2024limits}, afastando-se da superficialidade e da histeria tecnológica.

Mais do que um guia instrumental, o INATICON visa fomentar a discussão sobre a valorização docente em um cenário de crescente presença das IAs, conferindo maior segurança e liberdade aos educadores. Ao orientar o uso planejado da IA como um meio para ampliar o \textit{conatus} e a autonomia do professor, a pesquisa também se posiciona como uma ferramenta para a expansão dos multiletramentos docentes. Isso se dá ao promover uma autoria sobre o uso dessas tecnologias, incentivando práticas que podem ser tanto defensivas quanto vantajosas, ou, alternativamente, a abstenção consciente do uso em favor da exploração de outros meios de letramento em um ritmo que não seja acelerado.


\printbibliography\label{sec-bib}
% if the text is not in Portuguese, it might be necessary to use the code below instead to print the correct ABNT abbreviations [s.n.], [s.l.]
%\begin{portuguese}
%\printbibliography[title={Bibliography}]
%\end{portuguese}


%full list: conceptualization,datacuration,formalanalysis,funding,investigation,methodology,projadm,resources,software,supervision,validation,visualization,writing,review
\begin{contributors}[sec-contributors]
\authorcontribution{Leandro Paz da Silva}[conceptualization,datacuration,formalanalysis,investigation,methodology,software,validation,visualization,writing,review]
\authorcontribution{Márcia Amaral Corrêa Ughini Villarroel}[conceptualization,writing]
\authorcontribution{Fabio Yoshimitsu Okuyama}[methodology]
\end{contributors}

\section*{\fontsize{10}{12}\selectfont Agradecimentos}
Agradecimentos aos participantes da pesquisa anônimos. Agradecimentos aos organismos unicelulares e pluricelulares que têm resíduos da matéria livre e transformada do que um dia foi Benedictus deSpinoza em ação.

\begin{dataavailability}
\txtdataavailability{dataavailable} % options: dataavailable, dataonly, databody, datanotav, nodata
\end{dataavailability}


\end{document}

